\chapter{Design goals}

\label{ch:design-goals}

This chapter answers a single question: what is pii \emph{for}, and which constraints follow from
that purpose whether or not anyone likes them. Every later chapter describes a mechanism; this one
describes the pressure that shaped it. The purpose is narrow. pii exists so that a claim about work
a model did --- \enquote{the spec is implemented}, \enquote{the tests pass}, \enquote{the document
is current} --- can be checked by someone who was not present, from records the system kept while
it worked. Everything else in the workspace is downstream of that.

The book is deliberately organised around failure. Each of the five chapters that follow names the
mechanism \emph{and} the defect that made the mechanism necessary, because in this repository the
two are not separable: the exit rule in composition exists because a refusal silently vanished, the
health record exists because a status was being recomputed from 41{,}150 events on every call, and
the queue's lease exists because one dead claim wedged an entire lane. A design goal that has never
cost anybody anything is a slogan. The ones below have all been paid for.

\section{The four goals}

\begin{enumerate}
	\item \textbf{Evidence over narration.} A run leaves a journal of transitions; the stores keep
	counters and headers. A statement that cannot be reconstructed from those records is treated as
	unknown rather than believed.
	\item \textbf{Composition over configuration.} A graph is assembled from nodes and graphs that
	already exist. There is one implementation of \enquote{implement a spec} and several ways to
	reach it, never two implementations that can drift apart.
	\item \textbf{Refusal over guessing.} An unknown configuration key, an unserved model id, a
	port a node does not declare, a capability the graph does not hold, a pin that no longer matches
	the file: each refuses by name, at the earliest moment it can be detected.
	\item \textbf{Bounded growth.} Every store, journal, artefact directory and source file has an
	owner and a ceiling. An unattended host must not fill its disk, its memory or its review
	capacity while nobody is looking.
\end{enumerate}

The four are not independent. Evidence is what makes a refusal meaningful --- a refusal that leaves
no record is indistinguishable from a crash. Composition is what makes evidence cheap, because a
called graph journals under the caller's run id instead of inventing a second history. Bounded
growth is what keeps the first three affordable over months rather than days.

\begin{center}
\resizebox{\linewidth}{!}{%
\begin{tikzpicture}[node distance=8mm and 12mm]
	\node[b] (claim) {a claim\\\scriptsize \enquote{the spec is implemented}};
	\node[b, right=of claim] (run) {a run\\\scriptsize graph + input};
	\node[s, right=of run] (journal) {a journal\\\scriptsize one event per transition};
	\node[b, right=of journal] (replay) {a replay\\\scriptsize the same state, rebuilt};
	\node[v, below=of replay] (verdict) {agrees?};
	\node[b, left=of verdict] (done) {\texttt{done}\\\scriptsize checkable};
	\node[b, below=of verdict] (unknown) {\texttt{unverified}\\\scriptsize the record is silent};

	\draw[f] (claim) -- (run);
	\draw[f] (run) -- (journal);
	\draw[f] (journal) -- (replay);
	\draw[f] (replay) -- (verdict);
	\draw[f] (verdict) -- node[p,above=2pt,midway]{yes} (done);
	\draw[f] (verdict) -- node[p,right=2pt,midway]{no} (unknown);
	\draw[sf] (journal.south) to[bend right=22] node[p,below=2pt,midway]{rebuild} (replay.south);
\end{tikzpicture}}
\captionof{figure}{The chain a verdict has to survive: the dashed edge is what makes it a test rather than an
essertion, because the state is rebuilt from the journal and compared with the record the write path
maintained.}

\end{center}

\section{Evidence over narration}

A shell script leaves a transcript, and a transcript is not evidence. It says what was printed, not
which decision was taken at which port, and it cannot distinguish \enquote{the command succeeded}
from \enquote{the command was never run}. The engine therefore journals a \emph{state transition}
per node: the node, the port it left on, the value it produced, the hash of that value, and a
monotonic sequence number. From that record three things become possible that a transcript cannot
support: a run can be resumed from the position the journal says is next, a refusal can be
attributed to the hook that raised it, and health can be rebuilt rather than trusted.

The third of those is the one that was got wrong first, and the measurement is worth stating. Health
used to replay every journal of every run on every call. On this host, with 1{,}895 runs and
41{,}150 events in the store, that took 82, 85 and 92 milliseconds across three consecutive calls:
linear in history, on one of four async workers, with no cap. At ten times the history it is roughly
nine tenths of a second of a blocked worker per health call, and health is the surface every
watchdog reaches for first. The fix was not a cache with a time-to-live. It was the observation that
the answer is a \emph{status}, and a status should be maintained at the moment it changes, not
recomputed from first principles. The record is now written by the write path; the replay survives
as the test that the record still agrees with the events.

\begin{center}
\resizebox{\linewidth}{!}{%
\begin{tikzpicture}[node distance=8mm and 14mm]
	\node[b] (append) {event appended\\\scriptsize seq $n$ assigned};
	\node[b, right=of append] (status) {header status\\\scriptsize written in the same transaction};
	\node[b, right=of status] (read) {a reader\\\scriptsize one record};
	\node[s, below=of append] (journal) {journal\\\scriptsize append-only};
	\node[b, below=of status] (recount) {recount\\\scriptsize every event, once};
	\node[b, right=of recount] (compare) {compare\\\scriptsize equal, or a finding};

	\draw[f] (append) -- (status);
	\draw[f] (status) -- (read);
	\draw[f] (append) -- (journal);
	\draw[f] (journal) -- (recount);
	\draw[f] (recount) -- (compare);
	\draw[sf] (status) -- node[p,right=2pt,midway]{the test} (compare);
\end{tikzpicture}}
\captionof{figure}{Two readers of one truth --- the fast path reads a maintained record, the slow path recomputes from
the journal --- kept together because a counter with no recomputation behind it is an assertion with
extra steps.}

\end{center}

The record can be interrogated directly, which is the practical test of whether it is evidence at
all rather than a log:

\begin{lstlisting}[language=bash]
pii run trace <run_id>              # every transition, in sequence order
pii run explain <run_id>            # why the run ended as it did
pii graph lint graphs/queue.toml    # the same rules a run start applies, over one file
\end{lstlisting}

There is a second, quieter consequence. Because the journal is the primary artefact, a terminal
verdict must not be a summary of intentions. The engine determines a run's output by the last
finished node that has no outbound data edge on the port it took --- the node whose value nobody
consumed. If a final node left on a non-\texttt{ok} port carrying an error, the run is a failure
even though a graph author might have expected \enquote{it reached the end}. The record decides, not
the shape of the picture.

\section{Composition over configuration}

The second goal is stated as a rule about implementations, not about files: one behaviour, one
place. A graph that needs to implement a spec calls the graph that implements specs. It does not
carry a copy of it with different node ids, because a copy is a second implementation and the two
will diverge the first time only one of them is fixed.

This is not hypothetical in this repository. The development queue used to carry its own
\texttt{dispatch} coder and its own \texttt{judge}, doing inline what \texttt{graphs/spec.toml}
already did, and the two drifted in exactly the way the rule predicts: the queue's copy ran the
row's criteria while the spec graph ran the spec's, so a row could fail while the work it described
was green. Row 364 removed the copy and replaced it with one pinned \texttt{call}. The graph that
resulted is shorter than the one it replaced.

Composition also has a naming obligation that configuration does not. In \texttt{graphs/night.toml}
there was a node whose id was \texttt{sdd} whose kind was \texttt{call} with \texttt{graph =
"queue"}. An operator reading the night would conclude that SDD ran. It did not; the queue ran, and
the queue is row-by-row work rather than the SDD process. Row 373's requirement is not \enquote{add
an sdd graph} but \enquote{make the name true}: the node called \texttt{sdd} now calls
\texttt{sdd}. A composition whose names do not describe what is composed is worse than no
composition, because it is read by people who will not check.

\begin{center}
\resizebox{\linewidth}{!}{%
\begin{tikzpicture}[node distance=8mm and 12mm]
	\node[b] (night) {\texttt{night}\\\scriptsize tick, report};
	\node[b, right=of night] (sdd) {\texttt{sdd}\\\scriptsize contract, verify, docs};
	\node[b, right=of sdd] (tdd) {\texttt{tdd}\\\scriptsize write, red, implement, green};
	\node[b, below=of sdd] (docs) {\texttt{docs}\\\scriptsize one coder};
	\node[b, below=of night] (queue) {\texttt{queue}\\\scriptsize select, call, record};
	\node[b, right=of queue] (spec) {\texttt{spec}\\\scriptsize code, verify, reverify};

	\draw[f] (night) -- node[p,above=2pt,midway]{call} (sdd);
	\draw[f] (sdd) -- node[p,above=2pt,midway]{call} (tdd);
	\draw[f] (sdd) -- node[p,right=2pt,midway]{call} (docs);
	\draw[f] (queue) -- node[p,above=2pt,midway]{call} (spec);
	\draw[sf] (spec.north) to[bend right=18] node[p,right=2pt,midway]{call} (docs.south);
	\draw[bk] (night.south) -- node[p,left=2pt,midway]{no} (queue.east);
\end{tikzpicture}}
\captionof{figure}{What composes into what, and the edge that had to be removed --- the red dashed edge is the old
\texttt{sdd} node pointing at the queue, a composition that ran and meant something else.}

\end{center}

The cost of this goal is stated plainly because it is real: to know what \texttt{graphs/queue.toml}
does, a reader must also read \texttt{graphs/spec.toml}, and trust that the pin names the revision
they read. There is no single file that answers \enquote{what happens when the queue runs}. That is
the price of having one implementation instead of two, and Chapter~\ref{ch:composition} pays it in
detail.

\section{Refusal over guessing}

The third goal is the one that produces the most text in the repository, because every refusal
carries a machine-readable rule code. A refusal is not a log line that says \enquote{invalid graph};
it is a named rule, a named node, and --- where a capability is missing --- the name of the missing
capability. The reason is operational: at three in the morning the difference between
\texttt{graph.call.widen} and \texttt{V1} is the difference between reading one line of code and
reading the loader.

\begin{center}
\resizebox{\linewidth}{!}{%
\begin{tikzpicture}[node distance=8mm and 12mm]
	\node[b] (load) {graph text};
	\node[v, right=of load] (toml) {parses?};
	\node[v, right=of toml] (kinds) {keys known?\\\scriptsize kinds registered?};
	\node[v, right=of kinds] (ports) {ports declared?};
	\node[v, below=of ports] (caps) {caps within\\\scriptsize the ceiling?};
	\node[v, left=of caps] (pins) {pins match?};
	\node[b, left=of pins] (run) {the run starts};
	\node[b, below=of toml] (refused) {refused\\\scriptsize at load};

	\draw[f] (load) -- (toml);
	\draw[f] (toml) -- node[p,above=2pt,midway]{yes} (kinds);
	\draw[f] (kinds) -- node[p,above=2pt,midway]{yes} (ports);
	\draw[f] (ports) -- node[p,right=2pt,midway]{yes} (caps);
	\draw[f] (caps) -- node[p,above=2pt,midway]{yes} (pins);
	\draw[f] (pins) -- node[p,above=2pt,midway]{yes} (run);
	\draw[bk] (toml) -- node[p,left=2pt,midway]{no} (refused);
	\draw[bk] (kinds.south) to[bend left=10] node[p,below=2pt,midway]{no} (refused.east);
	\draw[bk] (ports.south) to[bend left=16] (refused.east);
	\draw[bk] (caps.south) to[bend right=8] node[p,below=2pt,midway]{no} (refused.east);
	\draw[bk] (pins.south) to[bend right=14] (refused.east);
\end{tikzpicture}}
\captionof{figure}{Refusal happens before execution and in a fixed order, so the message is the first thing actually
wrong rather than the last thing that failed, and nothing on the top row has run yet.}

\end{center}

\begin{longtable}{@{}p{0.30\textwidth}p{0.24\textwidth}p{0.36\textwidth}@{}}
	\toprule
	\textbf{Rule code} & \textbf{Raised by} & \textbf{What it names} \\
	\midrule
	\endhead
	\texttt{graph.call.widen} & \texttt{fragment/call.rs} & the call, the capability, the graph that
	does not hold it \\
	\texttt{graph.call.cycle} & \texttt{fragment/call.rs} & the chain of graphs, in order \\
	\texttt{graph.call.depth} & \texttt{fragment/call.rs} & the call and the depth cap of 4 \\
	\texttt{graph.call.exit} & \texttt{fragment/call.rs} & the callee with no \texttt{exit}, or an
	\texttt{exit} naming no node \\
	\texttt{graph.call.id} & \texttt{fragment/call.rs} & the node id the call would duplicate \\
	\texttt{graph.call.key} & \texttt{fragment/call.rs} & the one unknown key on the call \\
	\texttt{fragment.pin.mismatch} & \texttt{fragment/call.rs} & expected and actual sha256 \\
	\texttt{fragment.pin.drift} & \texttt{fragment/splice.rs} & the fragment path and both hashes \\
	\texttt{fragment.block.invalid} & \texttt{fragment/splice.rs} & a blocking fragment at a
	pre-surface \\
	\texttt{fragment.autonomy.forbidden} & \texttt{fragment/check.rs} & the node that would own a
	timer, store or budget \\
	\texttt{fragment.nesting.forbidden} & \texttt{fragment/check.rs} & the fragment that attaches or
	calls \\
	\texttt{graph.grant.ceiling} & \texttt{toml\_load.rs} & the graph or node, and the ceiling \\
	\texttt{V0}--\texttt{V11} & \texttt{validate.rs} & the rule, then the node \\
	\bottomrule
\end{longtable}

The codes share a property that matters more than their spelling: they are stable, greppable, and
each one maps to exactly one place in the source. That is what allows a test to assert a refusal by
name instead of asserting that \enquote{an error occurred}, which is the assertion that lets a
wrong refusal pass.

\section{Bounded growth}

The fourth goal is the least glamorous and the one that decides whether the other three survive a
year of unattended operation. Three ceilings are visible from the outside:

\begin{itemize}
	\item \textbf{Source files.} A file may have at most 450 effective lines --- non-blank,
	non-comment. The limit is enforced by a gate step, not by review, and the fix when it is hit is
	an extraction into a module that owns the concept. \texttt{daemon/sys\_status.rs},
	\texttt{store/sys\_status.rs} and \texttt{daemon/model\_roles.rs} all exist because a file
	crossed the line and the concept inside it deserved a name.
	\item \textbf{Artefacts.} A trajectory directory has an aggregate byte bound and a file-count
	bound, and the oldest finished runs are pruned past them. The defaults are 512 MiB and 50{,}000
	files; both are configuration keys rather than constants, because the right number depends on
	the host.
	\item \textbf{History.} Journals are append-only by design, so growth is managed by retention
	and compaction (\texttt{pii store compact}, \texttt{pii store prune}) rather than by rewriting
	events. An event that is deleted is a claim that can no longer be checked, which is why the
	operation is explicit and why the daemon must be stopped for it.
\end{itemize}

The ceilings are visible from the operator's side too, which is how a bound becomes a habit rather
than a rule:

\begin{lstlisting}[language=bash]
pii store stats                    # tables, sizes, and what is growing
pii store prune                    # retention, with the daemon stopped
bash quality-gate.sh --summary     # one line per gate step, including the file budget
\end{lstlisting}

Bounded growth is also why the engine has a step limit, a wall-clock limit and a token budget
rather than trusting nodes to be reasonable. A graph that loops is not an error in this system; a
graph that loops forever is, and the difference is a number in the graph's \texttt{[limits]} table.

\section{The cost, stated as a ledger}

Every goal above is bought with something. The table is the honest version of the trade, and each
row is developed in the chapter named in the last column.

\begin{longtable}{@{}p{0.22\textwidth}p{0.30\textwidth}p{0.24\textwidth}p{0.16\textwidth}@{}}
	\toprule
	\textbf{Choice} & \textbf{What it buys} & \textbf{What it costs} & \textbf{Chapter} \\
	\midrule
	\endhead
	Journal every transition & replay, attribution, rebuilt counters & write volume; a write-behind
	buffer that must be flushed at the right moments & \ref{ch:engine} \\
	Inline calls at load time & one implementation, no serialisation, one journal & a reader must
	follow a pin to know what runs; the callee becomes part of the caller's name space &
	\ref{ch:composition} \\
	Capabilities as a ceiling & an escalation must be written where it is reviewed & a graph that
	needs one more capability must be edited at two levels & \ref{ch:node-kinds} \\
	Named refusals & a failing test can assert the reason & more strings to keep stable than a
	single \enquote{invalid} & \ref{ch:node-kinds} \\
	450-line files & a concept has one home & more modules, more \texttt{mod} lines, more files to
	open & \ref{ch:workspace} \\
	Maintained counters & constant-time health & a write path that must never forget to update
	them; a rebuild test that must exist & \ref{ch:engine} \\
	\bottomrule
\end{longtable}

The last row is the one that recurs. Almost every defect described in this manual is a variant of
the same shape: a fast path added for a good reason, with no slow path left to check it. The
maintained counter, the inlined graph, the pinned fragment and the bounded file all trade a
computation for a record. The discipline that keeps the trade honest is that the computation is
kept somewhere as a test.

\section{The vocabulary, fixed early}

Five words carry most of the meaning in this book, and each one is used in a narrower sense than
its everyday meaning. They are stated once here so that later chapters can use them without
redefinition.

\begin{description}
	\item[node] A \texttt{kind} plus configuration. It does one thing, declares the ports it can
	leave from, and is the only unit the scheduler dispatches.
	\item[port] A named outcome, not a boolean. \texttt{select} has four; a coder has two; a judge
	has two or three. An edge names the port it leaves from.
	\item[graph] A table of nodes, a table of edges, an \texttt{entry}, a \texttt{caps} list,
	optional \texttt{exit} and optional \texttt{[limits]}. It is data, not code.
	\item[run] One execution of a graph with one input. It has an id, a journal, a header and
	exactly one terminal disposition.
	\item[call] A node kind that is replaced, at load time, by the nodes of another graph. It is not
	a subprocess and not a remote run.
	\item[fragment] A graph file under \texttt{graphs/fragments/} with a \texttt{[fragment]}
	header, attached to a host at a named surface and woven in at load time.
	\item[pin] The sha256 of a file's bytes, written into the file that includes it. A pin that
	does not match is a refusal, not a warning.
	\item[store] A redb database with named tables. Every store has a single writer and an owner
	crate.
\end{description}

Two of these definitions do quiet work. \emph{Port} being a name rather than a boolean is what makes
a graph read as a sentence: \texttt{red.ok} to \texttt{not-red} says something a \texttt{true}
branch cannot. And \emph{pin} being a hash of bytes rather than a version number is what makes a
review meaningful: a version can be reused, a byte string cannot change without the hash changing.

\section{Five defects that shaped the design}

The goals above are abstractions. The table below is the concrete version: five rows, each of which
changed a mechanism, with the symptom as an operator would have seen it.

\begin{longtable}{@{}p{0.06\textwidth}p{0.36\textwidth}p{0.50\textwidth}@{}}
	\toprule
	\textbf{Row} & \textbf{Symptom} & \textbf{What changed} \\
	\midrule
	\endhead
	353 & A claimed queue row never completed, and every other row stopped being claimed. A restart
	did not help, because a row that is \texttt{Running} blocks the queue by design. & The claim
	became a lease: a row is requeued when the run that claimed it is terminal or gone, and never
	while it is alive or suspended. A boot pass re-offers suspended runs without waiting for an
	event. \\
	364 & A row stayed \texttt{Running} although the called graph had refused. Nothing in the journal
	said why, because nothing after the refusal had run. & A host edge may name a callee's refusal
	port; the load marks that refusal node to return its verdict on \texttt{ok} instead of failing
	the run. \\
	370 & \texttt{pii health} took 82--92 ms and grew with history; at ten times the store it would
	block a worker for most of a second. & The totals became a record maintained by the write path;
	the replay stayed as the test that the record agrees with the events. \\
	372 & The engine could not express \enquote{this test must fail}, so a RED could not be proven
	and a test written after the code was indistinguishable from one written before it. & The judge
	accepts \texttt{expect = "fail"} and gains a third port, \texttt{unrunnable}, for a criterion
	that could not run at all. \\
	373 & A node named \texttt{sdd} called \texttt{queue}. The night claimed to run SDD and ran
	row-by-row work. & The name was made true: the node calls the \texttt{sdd} graph, and the pin
	moved in the same commit. \\
	\bottomrule
\end{longtable}

Read as a set, the five share a shape. In each case the system was doing something reasonable and
the \emph{record} was wrong: the row said \texttt{Running} when it was parked, health said nothing
about how it was obtained, the night's node said \texttt{sdd} when it was the queue. The fix is
never to make the code more clever; it is to make the record say what happened, and then to leave a
test that fails if the record and the events disagree.

\section{What pii is not}

Three negative statements prevent most misreadings of the rest of the book.

\begin{itemize}
	\item \textbf{It is not a workflow engine with an iteration primitive.} Composition is static:
	pieces are spliced at load time, so a graph cannot say \enquote{do this call once per criterion}.
	Everything that needs repetition today is one node looping internally --- a judge running four
	commands --- which is precisely why the judge has to know how to run criteria at all.
	\item \textbf{It is not a distributed system.} A call is inlined into the caller; there is no
	remote invocation, no queue between graphs and no serialisation boundary. The only subprocess in
	the picture is an agent, and the engine treats it as a node that may block, not as a peer.
	\item \textbf{It is not a scripting language.} A graph has no expressions beyond a small
	reference dialect (\texttt{\$.input.x}, \texttt{\$.nodes.n.field}) evaluated by the node that
	declares it. There are no local variables, no function calls and no conditionals outside a
	\texttt{route} or \texttt{gate} node. The restriction is the point: a graph that cannot compute
	can be checked by reading it.
\end{itemize}

\section{How the claims in this book were checked}

Every mechanism described here was read in the source and, wherever it could be, exercised. The
distinction matters because this repository has paid three times for a check that passed while the
live path was broken: a check that hand-wrote the lines it was meant to observe, a check fed a
source constant instead of the message the live path delivers, and a verification harness that was
itself wrong three times over. Each one certified a real defect as fixed.

The rule that follows from those three is stated in the workspace's agent guide and applies to
prose as much as to tests: a check's input must come from the live path, and every rule that skips
or suppresses something needs a near-miss case that must still take the normal path. The table
below is the version of that rule applied to the claims in this book.

\begin{center}
\footnotesize
\begin{tabular}{@{}p{0.30\textwidth}p{0.30\textwidth}p{0.30\textwidth}@{}}
\toprule
	\textbf{Claim} & \textbf{How it was checked} & \textbf{The near miss} \\
	\midrule
	A call returns through every node with an unwired \texttt{ok} port & read in
	\texttt{fragment/splice.rs}, then traced on \texttt{graphs/spec.toml}'s tail & a callee whose
	tail is one declared exit still returns on that node \\
	A wired refusal returns instead of failing the run & the unit tests in
	\texttt{fragment/call\_output.rs} assert the inner node gains \texttt{returns} & a refusal no host
	wires still fails the run \\
	A value naming a longer id is untouched & \texttt{rewrite\_str} refuses to match
	\texttt{\$.nodes.sddx} against \texttt{\$.nodes.sdd} & a bare \texttt{\$.nodes.sdd} is rewritten \\
	Health no longer scans journals & the requirement's test compares a rebuilt record with a replay
	& an absent record rebuilds rather than reporting zeros \\
	A judge can be expected to fail & ports change only when \texttt{spec} and \texttt{expect} are
	both present & a judge without \texttt{expect} keeps two ports \\
	\bottomrule
\end{tabular}
\end{center}

Where a claim could not be exercised --- the daemon's behaviour under a supervisor that kills it
mid-agent, for instance --- this book says so in the text rather than rounding it to a certainty.
The editor's rule for the volume is the same as the gate's: a number with no measurement behind it
is not allowed to look like one.

\section{What the reader should expect}

The chapters that follow are ordered by dependency rather than by importance. The workspace
(Chapter~\ref{ch:workspace}) explains which crate owns which behaviour and how to find a behaviour
from a symptom. The engine (Chapter~\ref{ch:engine}) explains what a run is: its states, its tick
loop, its journal, its recovery. Node kinds (Chapter~\ref{ch:node-kinds}) explain the vocabulary a
graph is written in --- what a node may declare and what it may do. Composition
(Chapter~\ref{ch:composition}) explains how graphs become each other's parts, and contains the one
rule in this system that surprises experienced readers: a call returns through every node that
declares an \texttt{ok} port and has no outgoing \texttt{ok} edge, which is not the same set as the
node the callee names in \texttt{exit}.

That rule is worth flagging here because it is the clearest example of the manual's method. It was
not designed; it was discovered, by a queue row that stayed \texttt{Running} for a night while the
graph it called had already refused. The next four chapters describe the system that produced that
behaviour, and the reasoning that fixed it without inventing a second inclusion mechanism.
